\section{Introduction}
\label{sec:introduction}

Density-based topology optimisation for Stokes flow, introduced by
Borrvall and Petersson~\cite{borrvall2003}, has been extended to
conjugate heat transfer~\cite{alexandersen2014}, scalar transport
problems~\cite{makhija2015}, and unsteady flows~\cite{kreissl2011}.
The Brinkman--SIMP approach replaces solid regions with a high-permeability
porous medium, enabling gradient-based optimisation of fluid topology
without remeshing.
Its application to microfluidics is particularly attractive because the
small length scales of soft-lithography devices make manual re-design
expensive and simulation-driven optimisation increasingly accessible~\cite{whitesides2006}.

Despite growing interest, several computational gaps remain in the
open-source literature on coupled flow-transport topology optimisation.

\subsection*{Gaps addressed by this work}

\emph{Gap~1 --- Open-source implementation of coupled
flow-transport topology optimisation.}
Existing open-source topology optimisation codes for fluidic systems
target pure flow~\cite{borrvall2003} or global scalar
mixing~\cite{na2022}.
FEniCSx-based frameworks exist for structural problems~\cite{dokken2023}
and Stokes flow~\cite{lazarov2016}, but a documented, tested,
and reproducible open-source implementation for the coupled
Brinkman--convection-diffusion system with a boundary
concentration-mismatch objective has not been published.
The linear gradient benchmark is chosen for verifiability
as a sanity check; the framework's primary contribution
is the implementation methodology, not the benchmark RMSE.

\emph{Gap~2 --- High-P\'{e}clet transport in FEniCSx TO.}
Existing FEniCSx topology optimisation
implementations~\cite{dokken2023} operate at $\Pe\ll1$
and do not require transport stabilisation.
At microfluidic P\'{e}clet numbers ($\Pe\sim10^2$--$10^5$),
SUPG stabilisation of both forward and adjoint equations
is required~\cite{brooks1982,donea2003}; we document
this standard technique in the FEniCSx UFL context.

\emph{Gap~3 --- End-to-end reproducibility.}
Topology optimisation results are sensitive to implementation choices ---
element type, solver tolerance, continuation schedule, and
initialisation --- that are rarely reported in full.
Published microfluidic topology optimisation results~\cite{na2022,feppon2024}
are not accompanied by open-source code, making independent
verification and extension difficult.
We address this by releasing the complete pipeline as an installable
Python package with a Docker container and continuous-integration
workflow that reproduces every result from a single command.

Table~\ref{tab:comparison} positions \texttt{micrograd} relative to existing frameworks.

\begin{table}[h]
\centering\footnotesize
\caption{Architectural comparison of open-source topology optimisation
         frameworks for microfluidic and flow-transport problems.
         \checkmark~=~implemented; {---}~=~not reported.
         Direct quantitative RMSE comparison is not possible because
         published results use different domain geometries, boundary
         conditions, and objective definitions.}
\label{tab:comparison}
\begin{tabular}{@{}p{2.4cm}p{1.4cm}p{1.4cm}p{1.4cm}p{1.4cm}p{1.4cm}@{}}
\toprule
\textbf{Feature} & BP03 & Na22 & Fe24 & dolfin- & \textbf{ours} \\
 & \cite{borrvall2003} & \cite{na2022} & \cite{feppon2024} & adj.~\cite{farrell2013} & \\
\midrule
Coupled flow-transport & --- & \checkmark & \checkmark & --- & \checkmark \\
Concentration-mismatch obj. & --- & \checkmark & --- & --- & \checkmark \\
SUPG adjoint stabilisation & --- & --- & --- & --- & \checkmark \\
Continuous adjoint & \checkmark & --- & \checkmark & --- & \checkmark \\
Discrete adjoint & --- & --- & --- & \checkmark & future work \\
FEniCSx implementation & --- & --- & --- & \checkmark & \checkmark \\
Docker + CI reproducibility & --- & --- & --- & --- & \checkmark \\
Open source & partial & --- & --- & \checkmark & \checkmark \\
\bottomrule
\end{tabular}
\end{table}

\subsection*{Contributions}

\begin{enumerate}
    \item \textbf{Continuous adjoint implementation for coupled
          Brinkman--convection-diffusion.}
          We derive and implement the continuous adjoint for the
          coupled system with a multi-objective cost functional,
          including the momentum--concentration coupling term
          $-\lambda\nabla c$ and a Dirichlet outlet condition
          derived from the misfit functional.
          While the adjoint structure follows standard Lagrangian
          methodology~\cite{borrvall2003,othmer2008}, a complete
          open-source FEniCSx implementation for this coupled system
          has not been previously published
          (Section~\ref{sec:adjoint}).

    \item \textbf{SUPG-stabilized adjoint implementation in FEniCSx.}
          We implement the stabilised forward and adjoint solvers
          and verify gradient accuracy via finite-difference Taylor
          test (slope $\approx 2.0$) and direction test
          (correlation $\geq 0.80$), confirming consistently implemented
          chain-rule propagation through Helmholtz filter and
          Heaviside projection
          (Section~\ref{sec:verification}).

    \item \textbf{Empirically validated continuation schedule.}
          A five-component schedule (Helmholtz filtering,
          Heaviside $\beta$-sharpening, $\alpha_{\max}$ ramping,
          hybrid OC/MMA, sinusoidal initialisation) with each
          component's motivation documented
          (Table~\ref{tab:continuation}).
          The hybrid OC/MMA split is validated by direct comparison:
          pure MMA gives RMSE~$=\rmsePureMMA$ vs.
          RMSE~$=\rmseShortRun$ for the hybrid
          (Table~\ref{tab:mma_comparison}).

    \item \textbf{Open-source reproducible implementation.}
          \texttt{micrograd} is released under the MIT licence
          with Docker, GitHub Actions CI, and Zenodo archival.
          All figures and results in this paper are recoverable
          from a single command in approximately one hour
          on a standard laptop
          (Section~\ref{sec:reproducibility}).
\end{enumerate}

\paragraph{Nature of contribution.}
This paper is a \emph{software implementation paper}.
The contribution is not a new optimisation algorithm,
not a record-breaking RMSE, and not a manufacturable device.
The contribution is a documented, tested, and reproducible
open-source implementation of adjoint-based topology optimisation
for a coupled Brinkman--convection-diffusion system in FEniCSx ---
something that does not currently exist in the open-source literature.

\paragraph{Scope and limitations.}
\texttt{micrograd}\footnote{The name \texttt{micrograd} coincidentally
matches an existing educational automatic-differentiation library
by Karpathy~\cite{karpathy2020micrograd};
this work is entirely independent and unrelated.}
operates within the Brinkman--SIMP porous-medium surrogate
on a fixed $80\times20$ mesh.
Two limitations are explicitly acknowledged and quantified:
(i)~the optimised permeability field is not a manufacturable design ---
hard-thresholding to binary permeability increases RMSE by $+\binaryGapPct\%$
because the porous optimum exploits gray-zone diffusion
unavailable in binary geometries; and
(ii)~the optimisation is not mesh-convergent ---
finer meshes find different local minima (RMSE $0.091$--$0.107$
on the $160\times40$ mesh vs.\ $\rmseTopOpt$ on the primary mesh).
These are known properties of density-based topology optimisation,
not defects of the implementation.
Both are documented, quantified, and identified as primary future directions
(robust projection~\cite{wang2011}, discrete adjoint~\cite{farrell2013}).

The remainder of this paper is organised as follows.
Section~\ref{sec:model} presents the mathematical model and adjoint
derivation.
Section~\ref{sec:methods} describes the FEniCSx implementation and
continuation strategies.
Section~\ref{sec:verification} presents verification studies.
Section~\ref{sec:results} presents the benchmark application.
Section~\ref{sec:discussion} discusses limitations and future directions,
and Section~\ref{sec:conclusion} concludes.